\documentclass[12pt]{article}
\usepackage[utf8]{inputenc}
\usepackage{graphicx}
\usepackage[numbers]{natbib}
\usepackage{amsmath}
\usepackage{amssymb}
\usepackage{mathrsfs}
\usepackage{hyperref}
\usepackage{natbib}
\usepackage{listings}
\usepackage{enumitem}
\usepackage{pgfplots}
\pgfplotsset{compat=1.18}
\usepackage{multirow}
\usepackage{booktabs}
\usepackage{algorithm}
\usepackage{algpseudocode}
\usepackage{tikz}
\usepackage{tabularx}
\usepackage{verbatim}
\usetikzlibrary{positioning,fit,backgrounds,arrows.meta,calc}

% Load packages that style.sty uses but doesn't load
\usepackage{geometry}
\usepackage{xcolor}
\usepackage{fancyhdr}
\usepackage{lastpage}

% Load style file (uses: geometry, xcolor, titlesec, titling, fancyhdr commands)
\usepackage{style}

% --- explicit layers so backgrounds never cover content ---
\pgfdeclarelayer{bg}
\pgfdeclarelayer{fg}
\pgfsetlayers{bg,main,fg}

% Define firstpage style (style.sty references it but Overleaf needs it defined explicitly)
\fancypagestyle{firstpage}{
  \fancyhf{}
  \fancyfoot[R]{\thepage/\pageref*{LastPage}}
}

% Math operators
\DeclareMathOperator*{\argmin}{arg\,min}
\DeclareMathOperator*{\argmax}{arg\,max}

% Code style
\lstset{
    basicstyle=\ttfamily\small,
    keywordstyle=\color{blue},
    commentstyle=\color{green!60!black},
    stringstyle=\color{red},
    numbers=left,
    numberstyle=\tiny,
    frame=single,
    breaklines=true,
    backgroundcolor=\color{gray!10}
}

% Slide-optimized settings
\titleformat{\section}
  {\Large\bfseries\color{TAMUmaroon}}
  {}{0pt}{}

\setlength{\parskip}{0.5em}
\setlength{\itemsep}{0.2em}
\pagestyle{fancy}

\title{\Large Research Update}
\author{Gaoming Lin \\
\vspace{10pt} Advisor: Dr. Byung-Jun Yoon}
\date{01/26/2026}

\begin{document}

\maketitle
\thispagestyle{firstpage}

\newpage
\section*{1. System Model, Latent Uncertainty, and Decision Objective}

\textbf{Network-coupled swing equation (IEEE-14, second-order Kuramoto):}

For each bus $i=1,\dots,N$,
\[
\dot{\theta}_i(t)=\omega_i(t)
\]
\[
M\,\dot{\omega}_i(t)=
P_{m,i}
-\sum_{j=1}^{N}B_{ij}\sin\!\big(\theta_i(t)-\theta_j(t)\big)
-(D+K)\omega_i(t)
+u^{\text{probe}}_{\xi,i}(t)
+u^{\text{ctrl}}_{\gamma,i}(t).
\]

\textbf{Index interpretation:}
\begin{itemize}
\item $i$ indexes the \emph{bus whose dynamics are being written} ($i=1,\dots,N$).
\item $j$ indexes all buses electrically coupled to $i$; $B_{ij}=0$ if no line exists.
\item The summation represents the total electrical power exchanged between bus $i$ and its neighbors.
\end{itemize}

\textbf{Latent parameters (fixed per episode):}
\[
\vartheta=(M,K)\sim p(\vartheta), \qquad \vartheta\in\mathbb{R}_+^2,
\]
where $M$ is the equivalent system inertia and $K$ is the aggregate fast frequency response (droop-like gain).

\textbf{Known / fixed:}
network coupling $B_{ij}$ (IEEE-14), damping $D$, nominal injections $P_{m,i}$.

\vspace{0.4em}
\textbf{Planning-level control decision (fast frequency response):}
\[
u^{\text{ctrl}}_{\gamma,i}(t)=\gamma\,g_i\,\omega_i(t),
\qquad \sum_i g_i = 1,\ g_i\ge0.
\]

\textbf{Security-constrained requirement:}
\[
\gamma^*(\vartheta)=
\min_{\gamma}
\quad \text{s.t.}\quad
\max_t|\dot f(t)|\le r_{\max},\;
\min_t f(t)\ge f_{\min}.
\]

\textbf{Goal:}
design probing experiments that reduce uncertainty in the
\emph{decision-relevant quantity} $\gamma^*(M,K)$, rather than estimating
$(M,K)$ for its own sake.

\textbf{References:}
Kundur (1994); Dörfler \& Bullo (2014).

% =========================================================
\newpage
\section*{2. Finite-Horizon Sequential Active Probing Experiment}

\textbf{Episode structure:}
\[
\vartheta\sim p(\vartheta)\ \text{once and fixed},\qquad t=1,\dots,T.
\]

At step $t$, the experimenter chooses a probing design
\[
\xi_t=(b_t,A_t,T_p)\in\Xi=\mathcal{B}\times\mathcal{A}\times\{T_p\}.
\]

\textbf{Active probing signal (IBR-style power injection):}
\[
u^{\text{probe}}_{\xi,i}(\tau)=
\begin{cases}
A_t\,s(\tau;T_p), & i=b_t,\\
0, & i\neq b_t,
\end{cases}
\quad
s(\tau;T_p)=\frac12\!\left(1-\cos\frac{2\pi\tau}{T_p}\right).
\]

\textbf{Physical interpretation:}
the probe is a small, smooth active-power pulse injected by an inverter-based
resource to excite inertial and primary-frequency dynamics while remaining
within operational limits.

\textbf{Observation generation:}
simulate the swing equation under $(\vartheta,\xi_t)$, measure frequency
response, compute observation $y_t$.

\textbf{Information available to the designer:}
\[
h_t=\{(\xi_1,y_1),\dots,(\xi_t,y_t)\},
\]
while $(M,K)$ remain unobserved.

\textbf{Reference:}
Peng et al. (2024); Chakraborty et al. (2022).

\section*{3. Active Probing Parameter Table (Fixed Design Choices)}

\begin{center}
\renewcommand{\arraystretch}{1.35}
\begin{tabular}{|c|c|c|p{4.6in}|}
\hline
\textbf{Parameter} & \textbf{Symbol} & \textbf{Value / Set} & \textbf{Justification (real references)} \\
\hline
Probe location &
$b_t$ &
$\mathcal{B}\subset\{1,\dots,14\}$ &
Buses with inverter-based actuation and high frequency observability, consistent with probing-based inertia estimation studies.
\newline
Chakraborty et al. (2022); Peng et al. (2024). \\
\hline
Probe amplitude &
$A_t$ &
$\mathcal{A}=\{A_1,A_2,A_3\}$ &
Selected so induced ROCOF exceeds PMU noise floor while remaining within security limits; validated via PHIL experiments.
\newline
Chakraborty et al. (2022); Peng et al. (2024). \\
\hline
Probe duration &
$T_p$ &
$2~\mathrm{s}$ &
Excites inertial and fast frequency response dynamics while avoiding slower secondary control. Peng et al. (2024). \\
\hline
Probe shape &
$s(t)$ &
Hann window &
Smooth, band-limited excitation widely used in probing and system identification.
\newline
Peng et al. (2024). \\
\hline
Sampling rate &
$f_s$ &
$12~\mathrm{Hz}$ &
Consistent with PMU frequency and ROCOF monitoring standards.
\newline
ENTSO-E (2020); NASPI (2021). \\
\hline
Observation window &
$T_{\text{obs}}$ &
$[0,10]~\mathrm{s}$ &
Captures ROCOF peak and early frequency transient before secondary control dominates.
\newline
ENTSO-E (2020); Wang et al. (2022). \\
\hline
\end{tabular}
\end{center}

% =========================================================
\newpage
\section*{4. Observation Model and Feature Extraction (ROCOF-only)}

\textbf{PMU-like frequency measurement:}
\[
\Delta f_i(t)=\frac{\omega_i(t)}{2\pi}, \qquad t=n\Delta t,\ \Delta t=1/f_s.
\]

\textbf{General feature map:}
\[
\phi(\Delta f)=
[\mathrm{ROCOF}_{\max}, f_{\min}, t_{\text{nadir}}, t_{\text{settle}}, \dots].
\]

\textbf{ROCOF-only observation (used in this work):}
\[
\widehat{\dot f}(n)=
\frac{\Delta f((n+1)\Delta t)-\Delta f(n\Delta t)}{\Delta t},
\qquad
y_t=\mathrm{ROCOF}_{\max}
=\max_{n\in\mathcal{W}_t}|\widehat{\dot f}(n)|.
\]

\textbf{Rationale:}
ROCOF is directly governed by inertia and fast frequency response in the swing equation and is the primary observable used in modern inertia monitoring and probing studies.

\textbf{References:}
ENTSO-E (2020); Wang et al. (2022); Chakraborty et al. (2022).

% =========================================================
\newpage
\section*{5. Likelihood Modeling and Why DAD (not iDAD)}

\textbf{Key clarification:}
the swing equation defines a \emph{deterministic state evolution}.
Uncertainty enters through measurement noise and unmodeled dynamics.
We therefore define an explicit likelihood at the \emph{measurement level}.

\textbf{Deterministic simulator-to-feature map:}
\[
\mu(\vartheta,\xi_t)
=
\mathrm{ROCOF}_{\max}\big(\Delta f(\cdot;\vartheta,\xi_t)\big).
\]

\textbf{Explicit measurement likelihood (offline training):}
\[
y_t = \mu(\vartheta,\xi_t) + \varepsilon_t,
\qquad
\varepsilon_t\sim\mathcal{N}(0,\sigma^2),
\]
\[
p(y_t\mid \vartheta,\xi_t)=
\mathcal{N}\!\left(\mu(\vartheta,\xi_t),\sigma^2\right).
\]

\textbf{Why this is valid:}
even though the swing equation itself has no closed-form likelihood,
the induced feature $y_t$ admits a tractable probabilistic noise model.

\textbf{Why DAD is selected:}
\begin{itemize}
\item DAD requires numerical evaluation of $\log p(y\mid\vartheta,\xi)$, not an analytic physics likelihood.
\item An explicit, testable likelihood exists at the feature level.
\item The iDAD framework explicitly states that DAD should be used whenever such a likelihood is available.
\end{itemize}

\textbf{References:}
Foster et al. (2021); Ivanova et al. (2021).

% =========================================================
\newpage
\section*{6. DAD Training, MOCU, and End-to-End Workflow}

\textbf{Design policy (finite-horizon, non-myopic):}
\[
\xi_t=\pi_\phi(h_{t-1}), \qquad t=1,\dots,T.
\]

\textbf{History update:}
\[
h_t=h_{t-1}\cup\{(\xi_t,y_t)\}.
\]

\textbf{Offline posterior (for evaluation only):}
\[
p(\vartheta\mid h_T)\propto
p(\vartheta)\prod_{t=1}^T p(y_t\mid \vartheta,\xi_t).
\]

\textbf{Mean Objective Cost of Uncertainty (MOCU):}
\[
\mathrm{MOCU}(h_T)=
\mathbb{E}_{\vartheta\sim p(\vartheta\mid h_T)}
\big[\gamma^*(\mathcal{A}_T)-\gamma^*(\vartheta)\big].
\]

\textbf{Training objective (decision-aware DAD):}
\[
\phi^*=\arg\min_\phi\ \mathbb{E}[\mathrm{MOCU}(h_T)].
\]

\textbf{Workflow summary:}
\begin{itemize}
\item Offline: simulate episodes, evaluate likelihood, train $\pi_\phi$ to minimize terminal MOCU.
\item Online: apply $\pi_\phi$ sequentially using observed history only.
\item Baselines are myopic; DAD optimizes the entire probing sequence.
\end{itemize}




\end{document}
